% Reconstructed text from V5 PDF page 47; non-native source
\par\smallskip\noindent\rule{\textwidth}{0.4pt}\par{\sffamily\footnotesize Original V5 page 47 -- recovered text}\par\begin{flushleft}\footnotesize\ttfamily
\par\smallskip
\par\smallskip
\noindent Diagnostics and Serving Infrastructure for MAS 47\par
\par\smallskip
\par\smallskip
\par\smallskip
\noindent 10 operate on single-trajectory analysis. MASPrism’s signal-extraction approach achieves com-\par
\noindent petitive accuracy with a 0.6B model and zero output tokens, at 6.69× lower latency than com-\par
\noindent parable methods. This trend suggests that the diagnostic layer’s demand on the infrastructure\par
\noindent layer is decreasing: newer methods require only logged traces, not re-execution access. The\par
\noindent implication for cross-layer design is significant — if future diagnostic methods universally op-\par
\noindent erate on logged traces stored independently of GPU memory, the Tier C evidence-loss scenario\par
\noindent becomes architecturally impossible.\par
\par\smallskip
\noindent Agent-aware infrastructure is emerging but single-objective. Nine of 22 infrastructure sys-\par
\noindent tems incorporate agent-level metadata into their management policies (ScaleSim, TokenCake,\par
\noindent KVFlow, KVCOMM, ForkKV, TokenDance, KAIROS, the policy-driven runtime, LRAgent).\par
\noindent However, all optimize for serving metrics (throughput, latency, cache hit rate) rather than di-\par
\noindent agnostic value. The agent-awareness is used to improve serving predictions, not to protect\par
\noindent diagnostic state. This creates an architectural asymmetry: the infrastructure layer is increas-\par
\noindent ingly capable of agent-level reasoning, but the objective function it optimizes does not include\par
\noindent diagnostic value.\par
\par\smallskip
\noindent The cross-layer conflict is narrower than assumed. The evidence-loss formalization (Eq. 1)\par
\noindent reveals that most diagnostic-serving interactions fall into Tier A (independent operation) or\par
\noindent Tier B (shared resource, reconstructable evidence). Tier C evidence-loss scenarios require non-\par
\noindent reconstructable state co-managed by the serving layer — a condition that current architectures\par
\noindent generally avoid by storing traces independently. Among 82 surveyed systems, we identified no\par
\noindent system where KV-cache eviction demonstrably destroyed diagnostic evidence in a production\par
\noindent deployment.\par
\par\smallskip
\noindent Step-level attribution remains an open challenge. Across the 12 failure-localization meth-\par
\noindent ods, step-level accuracy is substantially lower than agent-level accuracy. The best step-level re-\par
\noindent sults are approximately 52\% (TrajAudit RootSE and CHIEF on algorithm-generated scenarios),\par
\noindent compared to agent-level accuracies exceeding 70\%. Who\&When reports step-level accuracy\par
\noindent degrades to approximately 0\% on trajectories exceeding 93 steps. TraceElephant demonstrates\par
\noindent that complete traces improve step-level accuracy by 76\% relative (16\% to 28.1\%), compared\par
\noindent to only 22\% relative improvement at the agent level. This granularity gap implies that in-\par
\noindent frastructure decisions affecting trace completeness have a disproportionate impact on the most\par
\noindent challenging — and often most useful — level of diagnostic attribution.\par
\par\smallskip
\noindent Benchmark infrastructure is fragmented. The 12 surveyed benchmarks collectively cover\par
\noindent failure attribution, reliability, and safety evaluation, but use different trace formats, annotation\par
\noindent protocols, evaluation metrics, and failure definitions. No common evaluation protocol exists,\par
\noindent and 0 of 12 benchmarks evaluate diagnostic-infrastructure interaction. TraceElephant [6] pro-\par
\noindent vides the richest traces (380 structured JSONs with step-level tool logs) but no serving-side\par
\noindent instrumentation. Who\&When [48] defines the first standardized attribution protocol but oper-\par
\noindent ates entirely offline. TRAIL [10] is the first OTel-format benchmark, enabling integration with\par
\noindent serving-side telemetry, but its best model achieves only 11\% joint accuracy — suggesting that\par
\noindent structured traces alone are insufficient without specialized attribution methods. The absence of a\par
\noindent cross-layer benchmark is both a gap and an opportunity: building one would require combining\par
\noindent a realistic multi-agent workload generator, a configurable serving backend (with tunable evic-\par
\noindent tion policy and memory pressure), injected failures at known locations, and matched diagnostic\par
\noindent and serving metrics.\par
\end{flushleft}
